\documentclass[11pt]{article}
\usepackage[margin=1in]{geometry}
\usepackage{amsmath,amssymb,amsthm}
\usepackage{graphicx}
\usepackage[hidelinks]{hyperref}
\usepackage{microtype}

\newtheorem{theorem}{Theorem}
\newtheorem{lemma}{Lemma}
\newtheorem{corollary}{Corollary}

\newcommand{\R}{\mathbb{R}}
\newcommand{\e}{\mathfrak e}

\title{Sharp Extremal Gap Counts at Fixed Essential Period\\
\large A partial answer to Section 5.1 of arXiv:2112.15055}
\author{Autonomous proof packet for \texttt{fa\_banach\_001}}
\date{17 August 2026}

\begin{document}
\maketitle

\noindent\textbf{Status.}
Candidate partial result, likely valid, pending human review.  This packet
proves a sharp extremal theorem for the source's gap-versus-period question;
it does not classify every intermediate gap pattern.

\section{Source question}

Kumar and Kumar ask for a relationship between the number of spectral gaps
and the essential period of a self-adjoint periodic Jacobi operator and write
that the number of gaps may be expected to be ``on the order of'' the period
\cite[Section 5.1, PDF p.~21]{KK}.

\begin{center}
\includegraphics[width=0.98\textwidth]{figures/open_problem_crop.png}
\end{center}

For a two-sided self-adjoint Jacobi operator
\[
 (Ju)_n=a_{n-1}u_{n-1}+b_nu_n+a_nu_{n+1},
 \qquad a_n>0,\quad b_n\in\R,
\]
the \emph{essential period} is the least positive period of both coefficient
sequences.  An open spectral gap is a bounded component of the complement of
the spectrum between two adjacent spectral bands.

\section{Idea of the proof}

The period is encoded not by the number of bands, but by their equilibrium
(harmonic) masses.  A two-band set can have one band of mass exactly $1/p$.
Finite-gap inverse spectral theory then produces a periodic Jacobi matrix with
that set as spectrum, and the reduced denominator $p$ of the harmonic mass
forces its least period to be exactly $p$.  Thus one open gap survives at
arbitrarily large periods.  B\"org's theorem rules out zero gaps when the
least period exceeds one, so one is the sharp minimum.  Standard Floquet
theory gives at most $p-1$ gaps and the source exhibits an example attaining
that bound, giving the sharp maximum.

\section{Finite-gap input}

Let
\[
 \e=[\alpha_1,\beta_1]\cup\cdots\cup
 [\alpha_{\ell+1},\beta_{\ell+1}]
\]
be a finite-gap compact set, and let $\rho_{\e}$ be its logarithmic
equilibrium measure.  The quantities
$\rho_{\e}([\alpha_j,\beta_j])$ are its band harmonic measures.

We use the following standard consequence of the isospectral-torus
construction of Christiansen, Simon, and Zinchenko \cite[Corollaries 4.5 and
6.4]{CSZ}.

\begin{lemma}[finite-gap period criterion]\label{lem:fg}
For every finite-gap set $\e$ there are two-sided self-adjoint Jacobi matrices
in its isospectral torus with spectrum $\e$.  The coefficient sequences of
every member of the torus are periodic with period $q$ if and only if all band
harmonic measures belong to $q^{-1}\mathbb Z$.  Consequently their least
period is the least common denominator of those masses.
\end{lemma}

The last sentence follows directly from the ``if and only if'' assertion:
period $q$ is possible precisely when multiplication by $q$ makes every mass
integral.

\section{One gap at every least period}

\begin{lemma}\label{lem:mass}
For every integer $p\geq2$, there is a union $\e_p$ of two disjoint compact
intervals whose band equilibrium masses are $1/p$ and $(p-1)/p$.
\end{lemma}

\begin{proof}
For $p=2$, take any symmetric two-band set
\[
 [-1,-a]\cup[a,1],\qquad 0<a<1.
\]
Reflection symmetry and uniqueness of the equilibrium measure give mass
$1/2$ to each band.

Now let $p>2$ and, for $0<s<1$, put
\[
 E_s=[-1,0]\cup[s,1],\qquad
 h(s)=\rho_{E_s}([s,1]).
\]
The equilibrium measures vary weakly continuously when these interval
endpoints vary.  One quick justification is the standard energy-minimization
argument: along a monotone sequence of endpoints, logarithmic capacity is
continuous and every weak limit of the equilibrium measures has the minimal
energy of the limiting compact set; uniqueness identifies that limit.  An
arbitrary convergent sequence has monotone subsequences, giving full
continuity.  Since the two bands remain separated near an interior parameter,
their masses are continuous as well.

As $s\downarrow0$, the compact sets tend to $[-1,1]$ and
\[
 h(s)\longrightarrow \rho_{[-1,1]}([0,1])=\frac12
\]
by reflection symmetry.  As $s\uparrow1$, the right interval collapses to the
polar singleton $\{1\}$; the limiting equilibrium measure is that of
$[-1,0]$ and gives this singleton mass zero.  Hence $h(s)\to0$.  The
intermediate value theorem supplies $s_p\in(0,1)$ with $h(s_p)=1/p$.
Taking $\e_p=E_{s_p}$ proves the claim.
\end{proof}

\begin{theorem}[one-gap realization]\label{thm:onegap}
For every integer $p\geq2$, there is a self-adjoint Jacobi operator of least
period exactly $p$ whose spectrum has exactly one open gap.
\end{theorem}

\begin{proof}
Apply Lemma~\ref{lem:fg} to the two-band set $\e_p$ from
Lemma~\ref{lem:mass}.  Its band masses lie in $p^{-1}\mathbb Z$, so the
isospectral torus gives a $p$-periodic Jacobi matrix $J_p$ with
$\sigma(J_p)=\e_p$.

Its least period cannot be a positive integer $q<p$.  Indeed, period $q$
would force $1/p=k/q$ for an integer $k$, by Lemma~\ref{lem:fg}; equivalently
$q=kp$, impossible for $0<q<p$.  Thus the least period is $p$.  Since
$\e_p$ consists of exactly two disjoint bands, it has exactly one open gap.
\end{proof}

\section{Sharp extremal theorem}

For $p\geq2$, let $g_{\min}(p)$ and $g_{\max}(p)$ denote respectively the
minimum and maximum number of open gaps among self-adjoint Jacobi operators
of least period $p$.

\begin{theorem}\label{thm:extremal}
For every integer $p\geq2$,
\[
 \boxed{\quad g_{\min}(p)=1,\qquad g_{\max}(p)=p-1.\quad}
\]
\end{theorem}

\begin{proof}
B\"org's theorem for periodic Jacobi matrices says that all gaps are closed
only when the coefficient sequences are constant \cite{KK}.  A matrix of
least period $p>1$ therefore has at least one open gap.  Theorem
\ref{thm:onegap} attains one, proving $g_{\min}(p)=1$.

Floquet theory represents the spectrum of a $p$-periodic Jacobi matrix as at
most $p$ compact bands, so there are at most $p-1$ open gaps.  The explicit
matrix-valued symbol $f_0(\theta)$ in Section 5.1 of the source has least
period $p$ and exactly $p-1$ gaps \cite[p.~21]{KK}.  Hence the upper bound is
attained and $g_{\max}(p)=p-1$.
\end{proof}

\begin{corollary}
No estimate of the form
\[
 \#\{\text{open gaps of }J\}\geq F(p)
\]
can hold for all self-adjoint Jacobi matrices of least period $p$ if
$F(p)>1$ for any $p\geq2$.  In particular, no positive linear lower bound in
$p$ is possible.
\end{corollary}

\section{Verification, scope, and novelty}

The proof is analytic and uses no numerical computation.  The decisive checks
for a reviewer are (i) weak continuity and the endpoint limits of $h(s)$ in
Lemma~\ref{lem:mass}, and (ii) the exact-period inference from Corollaries 4.5
and 6.4 of \cite{CSZ}.  The source PDF and the decisive supporting PDF are
included locally.

The result determines the sharp minimum and maximum number of gaps at each
least period.  It does not classify all intermediate counts or band
geometries, and it does not address the source's separate non-normal
pseudospectral inequalities.

A bounded run-index and arXiv/web search on 17 August 2026 used the exact and
near-exact phrases ``one gap $p$-periodic Jacobi'', ``exactly one gap periodic
Jacobi'', ``minimal period one gap Jacobi matrix'', and ``arbitrary period
finite gap Jacobi''.  No exact statement of Theorem~\ref{thm:extremal} was
found.  The argument is an agent-identified synthesis of established
finite-gap theory, so novelty confidence is moderate.

\section*{Human-review recommendation}

Review is recommended.  If the equilibrium-measure continuity lemma and the
least-denominator reading of the finite-gap period criterion are confirmed,
the packet is a complete sharp extremal answer to the quantitative
gap-versus-period interpretation of the source question.

\begin{thebibliography}{9}
\bibitem{KK}
Krishna Kumar G. and V. B. Kiran Kumar,
\emph{Stability in Non-Normal Periodic Jacobi Operators: Advancing B\"org's
Theorem}, arXiv:2112.15055, revised 2024.

\bibitem{CSZ}
Jacob S. Christiansen, Barry Simon, and Maxim Zinchenko,
\emph{Finite Gap Jacobi Matrices, I. The Isospectral Torus},
arXiv:0810.3273, 2008.
\end{thebibliography}

\end{document}
